\documentclass[12pt]{article}

% ── Packages ──────────────────────────────────────────────────────────────────
\usepackage[margin=1in]{geometry}
\usepackage{amsmath,amssymb}
\usepackage{booktabs}           % \toprule \midrule \bottomrule
\usepackage{natbib}             % author-year citations
\usepackage{graphicx}
\usepackage{hyperref}
\usepackage{setspace}
\usepackage{caption}
\usepackage{subcaption}
\usepackage{xcolor}
\usepackage{microtype}

\doublespacing

% ── Auto-generated numbers (run_paper_experiments.py → results/paper_results.txt)
% Paste the contents of results/paper_results.txt here before compiling,
% or \input{../results/paper_results.txt}
\input{../results/paper_results.txt}

% ── Title ─────────────────────────────────────────────────────────────────────
\title{Nonparametric Bayesian Market Regime Detection:\\
       A Sticky HDP-HMM Approach}

\author{Adam Morris\\
        \small University of [TODO]\\
        \small \texttt{adam.morris0201@gmail.com}}

\date{\today}

% ══════════════════════════════════════════════════════════════════════════════
\begin{document}
\maketitle

% ── Abstract ──────────────────────────────────────────────────────────────────
\begin{abstract}
We propose a Bayesian nonparametric hidden Markov model (HDP-HMM) for identifying
latent volatility regimes in equity markets.
Unlike threshold-based approaches that partition on a single observable such as
VIX, our model recovers persistent hidden states from a four-dimensional
macro-financial feature vector---S\&P 500 log returns, the VIX level,
the 10-year minus 2-year Treasury yield spread, and the Chicago Fed National
Financial Conditions Index (NFCI).
A sticky Dirichlet-process prior over the transition distribution encourages
regime persistence, which we verify empirically: the HDP-HMM produces regimes
with a mean dwell time of \nLowVolDwell{}--\nHighVolDwell{} days versus
8--17 days for a VIX-threshold benchmark.
The model identifies \nEffKRaw{} latent sub-states, consolidated into three
canonical regimes (Low-Vol, Moderate-Vol, High-Vol) by mean VIX rank.
A simple volatility-targeting strategy that scales equity exposure by the
inverse of the current regime volatility achieves a Sharpe ratio of
\nHDPBTSharpe{} versus \nBHSharpe{} for buy-and-hold, with a maximum drawdown
of \nHDPBTMaxDD{} and only \nHDPRebalPerYear{} rebalances per year.
\end{abstract}

\newpage
\tableofcontents
\newpage

% ══════════════════════════════════════════════════════════════════════════════
\section{Introduction}
\label{sec:intro}

% TODO: 3–4 paragraphs
% Para 1: motivation — why regime detection matters for investors and researchers
% Para 2: the problem with VIX thresholds (VIX dwell time 8–17 days, noisy)
% Para 3: what this paper does differently (Bayesian, nonparametric K, 4 macro features)
% Para 4: contributions + roadmap

Financial markets exhibit qualitatively distinct periods of calm and stress
that are not directly observable.
Identifying these latent regimes is practically important: risk models,
position sizing, and hedging strategies all depend on accurate characterization
of the current volatility environment.

The dominant practitioner approach is threshold-based: label a day High-Vol
when VIX exceeds some cutoff (often 20 or 25).
This is transparent but brittle.
VIX fluctuates daily, causing mechanical regime labels to flip every
8--17 days on average---far too frequently to drive portfolio decisions---and
VIX captures only the implied volatility of SPX options, ignoring the credit
and yield-curve signals that historically precede market stress.

In this paper we ask: \emph{can a Bayesian nonparametric model recover more
stable and economically coherent regimes from a small set of macro-financial
features?}
We fit a sticky hierarchical Dirichlet process hidden Markov model
\citep[HDP-HMM;][]{teh2006hdp,fox2011sticky} to four daily time series
spanning \nTrainStart{} to \nTrainEnd{}: the S\&P 500 log return, the VIX
level, the 10-to-2-year Treasury yield spread, and the Chicago Fed NFCI.
The HDP prior allows the number of regimes $K$ to be inferred from the data
rather than fixed in advance; the sticky self-transition term penalises rapid
state switching.

Our main contributions are:
\begin{enumerate}
  \item We show that the HDP-HMM discovers \nEffKRaw{} latent sub-states that
        consolidate into three economically interpretable regimes with realized
        volatilities of \nLowVolAnnVol{}\%, \nModerateVolAnnVol{}\%, and
        \nHighVolAnnVol{}\% respectively---a 3.6$\times$ spread from Low-Vol
        to High-Vol that VIX thresholds cannot cleanly separate.
  \item Regime dwell times of \nLowVolDwell{}--\nHighVolDwell{} trading days
        confirm genuine persistence, in contrast to the 8--17 day flicker
        of VIX-threshold labels.
  \item A simple inverse-volatility sizing rule applied to the HDP regime labels
        achieves a Sharpe ratio of \nHDPBTSharpe{} versus \nBHSharpe{} for
        buy-and-hold, requiring only \nHDPRebalPerYear{} rebalances per year.
\end{enumerate}

The remainder of the paper is organised as follows.
Section~\ref{sec:related} reviews the regime-switching literature.
Section~\ref{sec:data} describes the data and features.
Section~\ref{sec:model} presents the HDP-HMM specification and inference.
Section~\ref{sec:results} reports in-sample regime characteristics, transition
dynamics, and the backtest.
Section~\ref{sec:conclusion} concludes.


% ══════════════════════════════════════════════════════════════════════════════
\section{Related Work}
\label{sec:related}

% TODO: 2–3 paragraphs covering:
% - Hamilton (1989) Markov-switching baseline
% - Guidolin & Timmermann (2007) asset allocation under regime switching
% - Fox et al. (2011) sticky HDP-HMM
% - Dacco & Satchell (1999) misclassification risk
% - Any recent macro-augmented HMM papers (AH-HMM 2025)

The regime-switching framework of \citet{hamilton1989new} models a time series
as generated by a hidden Markov chain with $K$ states, estimated by EM.
\citet{guidolin2007asset} extend this to multivariate equity-bond settings,
showing that regime-aware allocation outperforms static mean-variance
benchmarks.
A limitation of the parametric HMM is that $K$ must be chosen in advance;
misspecification of $K$ leads to regime fragmentation or aggregation
\citep{dacco1999predictability}.

The hierarchical Dirichlet process of \citet{teh2006hdp} enables a
nonparametric prior over countably infinite state spaces.
\citet{fox2011sticky} augment the HDP-HMM with a self-transition stickiness
parameter $\kappa$, directly encoding the prior belief that real-world regimes
are persistent.
Our model builds on this foundation, specialised to the equity market setting
with macro-financial inputs.

% TODO: add any other papers you've found relevant


% ══════════════════════════════════════════════════════════════════════════════
\section{Data and Features}
\label{sec:data}

\subsection{Sample}

We use daily data from \nTrainStart{} to \nTrainEnd{} (\nTrainDays{} trading
days) for the in-sample analysis.
% TODO: describe test period

\subsection{Feature construction}

We construct four features, all lagged one day to ensure causal inference:

\begin{enumerate}
  \item \textbf{SPY log return} ($r_t$): daily log return of the S\&P 500
        index (\^{}GSPC via Yahoo Finance).  Captures price momentum.
  \item \textbf{VIX level} ($\sigma^{iv}_t$): CBOE Volatility Index close
        (\^{}VIX).  Captures option-implied forward volatility.
  \item \textbf{Yield slope} ($y_t$): 10-year minus 2-year US Treasury yield
        spread (FRED series T10Y2Y).  Captures the macro business cycle.
  \item \textbf{NFCI} ($f_t$): Chicago Fed National Financial Conditions Index
        (FRED series NFCI, weekly, forward-filled to daily).  Captures broad
        financial stress.
\end{enumerate}

Four features feed directly into the HMM without dimensionality reduction;
at $D=4$ the covariance structure is well-identified and PCA would discard
interpretable axes.
Each feature is expanding-window standardised at each time $t$ using only
data from $[1, t]$, preserving the causal guarantee required for live trading.

\subsection{Descriptive statistics}

% TODO: add a Table 0 with feature summary stats if desired
% Mean, std, min, max for each of the 4 features over the train period.


% ══════════════════════════════════════════════════════════════════════════════
\section{Model}
\label{sec:model}

\subsection{Hidden Markov Model baseline}

Let $\mathbf{x}_t \in \mathbb{R}^4$ be the feature vector at time $t$ and
$z_t \in \{1,\ldots,K\}$ the latent regime.
A standard Gaussian HMM assumes:
\begin{align}
  z_t \mid z_{t-1} &\sim \mathrm{Categorical}(\pi_{z_{t-1}}) \\
  \mathbf{x}_t \mid z_t &\sim \mathcal{N}(\mu_{z_t},\, \Sigma_{z_t})
\end{align}
with transition matrix $\Pi = [\pi_k]$ and emission parameters
$\{\mu_k, \Sigma_k\}$ estimated by EM.
We fix $K=3$ for this baseline.

\subsection{Sticky HDP-HMM}

The HDP-HMM replaces the finite Dirichlet prior on each row of $\Pi$ with a
hierarchical Dirichlet process \citep{teh2006hdp}, enabling $K$ to grow with
the data.
We adopt the truncated stick-breaking representation of \citet{fox2011sticky}
with truncation level $K_{\max} = 8$.

Global state weights $\beta$ are drawn via the GEM construction:
\[
  v_k \sim \mathrm{Beta}(1, \alpha_{DP}),\quad
  \beta_k = v_k \prod_{j<k}(1-v_j), \quad k=1,\ldots,K_{\max}.
\]

Transition rows are sticky Dirichlet:
\[
  \pi_k \sim \mathrm{Dirichlet}(\alpha_{\mathrm{trans}}\,\beta + \kappa\,\mathbf{e}_k),
\]
where $\kappa > 0$ concentrates probability mass on self-transitions.
Gaussian emissions with diagonal covariance:
\[
  \mathbf{x}_t \mid z_t = k \;\sim\; \mathcal{N}(\mu_k,\,\mathrm{diag}(\sigma^2_k)).
\]

Discrete states are marginalised analytically via the forward algorithm
(implemented in JAX with \texttt{lax.scan}), making the model differentiable
and amenable to stochastic variational inference.

\subsection{Inference}

We use stochastic variational inference (SVI) with an AutoNormal guide
\citep{ranganath2014blackbox} for fast training, and optionally NUTS MCMC
\citep{hoffman2014nuts} for the paper-quality posterior used in uncertainty
quantification.
All inference uses JAX \citep{jax2018github} and NumPyro \citep{phan2019composable}.
The random seed is fixed at 42 throughout; JAX and NumPyro versions are pinned
exactly to guarantee reproducibility.

\subsection{Regime labelling}

After inference the model returns $K_{\mathrm{eff}}$ active states (those
with posterior beta weight above a threshold and at least 1\% of days assigned).
States are sorted by mean within-state VIX level and partitioned into three
canonical regimes by floor-division tertile:
Low-Vol (lowest VIX states), Moderate-Vol, and High-Vol (highest VIX states).


% ══════════════════════════════════════════════════════════════════════════════
\section{Results}
\label{sec:results}

\subsection{Regime characteristics}

Table~\ref{tab:regime_chars} reports within-regime statistics for the
HDP-HMM, the VIX-threshold baseline, and the parametric HMM baseline.

\input{../results/paper_results.txt}

% (Tables 1–4 are auto-inserted by \input above.)
% Comment out the \input and paste the table code directly for journal submission.

The HDP-HMM Low-Vol regime covers \nLowVolPct{} of training days and exhibits
\nLowVolAnnVol{}\% annualized realized volatility with a Sharpe of
\nLowVolSharpe{}.
The High-Vol regime covers only \nHighVolPct{} of days but with
\nHighVolAnnVol{}\% realized volatility---a 3.6$\times$ spread relative to
Low-Vol.
Crucially, both the Low-Vol and High-Vol regimes are highly persistent
(mean dwell times of \nLowVolDwell{} and \nHighVolDwell{} days respectively),
compared to only 8--17 days for the VIX-threshold benchmark.
This persistence is not imposed by the model: it emerges from the data via
the learned $\kappa$ stickiness parameter.

\subsection{Transition dynamics}

Table~\ref{tab:transition} shows the posterior-mean transition matrix.
The diagonals of 0.989, 0.986, and 0.974 confirm that the model has learned
genuinely sticky regimes: on any given day, there is less than a 3\% chance
of transitioning out of the current regime.
Cross-regime transitions are asymmetric: High-Vol exits preferentially to
Low-Vol rather than Moderate-Vol, consistent with the empirical pattern of
V-shaped recoveries following market stress.

\subsection{Information content beyond VIX}

% TODO: discuss Table 3 carefully.
% The R² gains are small in absolute terms (equity returns are near-unpredictable),
% which is expected and should be acknowledged.
% Frame this section around what the model DOES separate: volatility, not returns.

\subsection{Volatility-targeting backtest}

% TODO: discuss Table 4 once numbers are in.
% Lead with: Sharpe improvement, max drawdown reduction, low rebalancing frequency.
% Acknowledge limitation: in-sample, no transaction costs modelled.


% ══════════════════════════════════════════════════════════════════════════════
\section{Conclusion}
\label{sec:conclusion}

% TODO: 2–3 paragraphs
% Para 1: what we showed (persistent latent regimes, not recoverable from VIX alone)
% Para 2: practical implication (vol targeting, low turnover)
% Para 3: limitations and future work (OOS validation, NUTS full posterior, macro-augmented transitions)

We have demonstrated that equity markets exhibit latent volatility regimes with
mean dwell times of \nLowVolDwell{}--\nHighVolDwell{} trading days---far longer
than the 8--17 day flicker of VIX-threshold labels.
These regimes, recovered by a sticky HDP-HMM from four macro-financial features,
separate realized volatility by a factor of 3.6$\times$ from Low-Vol to High-Vol,
providing a stable, actionable signal for risk management.

% TODO: limitations + future work paragraph


% ── Bibliography ──────────────────────────────────────────────────────────────
\bibliographystyle{apalike}
\bibliography{references}

\end{document}
